\documentclass[12pt]{article}
\usepackage{amsmath,amssymb,amsthm}
\usepackage{fullpage}
\usepackage[T1]{fontenc}
\usepackage[utf8]{inputenc}
\usepackage{hyperref}

\newtheorem{theorem}{Theorem}
\newtheorem{lemma}[theorem]{Lemma}
\newtheorem{corollary}[theorem]{Corollary}
\newtheorem{proposition}[theorem]{Proposition}
\theoremstyle{definition}
\newtheorem{definition}[theorem]{Definition}
\theoremstyle{remark}
\newtheorem{remark}[theorem]{Remark}

\title{Solution to Problem 1:\\
Equivalence of the $\Phi^4_3$ Measure under Smooth Shifts}
\author{}
\date{April 2026}

\begin{document}
\maketitle

\begin{abstract}
We prove that the $\Phi^4_3$ measure $\mu$ on $\mathcal{D}'(\mathbb{T}^3)$ is
equivalent to its pushforward $T_\psi^*\mu$ under the shift $T_\psi(u) = u+\psi$
for any smooth $\psi:\mathbb{T}^3\to\mathbb{R}$.  The argument proceeds via a
chain of measure equivalences: $\mu\sim\mu_0\sim T_\psi^*\mu_0\sim T_\psi^*\mu$,
where $\mu_0$ is the Gaussian Free Field.  The key ingredients are the
Cameron--Martin theorem for $\mu_0$ and the positivity of the $\Phi^4_3$ density.
\end{abstract}

%------------------------------------------------------------------
\section{Setup and Statement}

Let $\mathbb{T}^3 = (\mathbb{R}/\mathbb{Z})^3$ be the three-dimensional unit torus.
Fix a mass $m>0$ and let $C = (-\Delta + m^2)^{-1}$ be the covariance operator of
the massive free field on $\mathbb{T}^3$.  The \emph{Gaussian Free Field} (GFF)
$\mu_0$ is the centred Gaussian measure on $\mathcal{D}'(\mathbb{T}^3)$ with
covariance $C$.

The \emph{$\Phi^4_3$ measure} is (formally)
\[
  d\mu(\phi) \;=\; Z^{-1}\exp\!\left(-\int_{\mathbb{T}^3} :\!\phi^4\!:(x)\,dx\right)d\mu_0(\phi),
\]
where $:\!\phi^4\!:$ denotes the Wick-renormalized fourth power of $\phi$ with
respect to $\mu_0$.  The construction of $\mu$ as a genuine probability measure
on $\mathcal{D}'(\mathbb{T}^3)$ is classical (Nelson, Glimm--Jaffe, Hairer).
We take as given the following well-known properties of $\mu$:

\begin{enumerate}
\item[\textbf{(P1)}] (\emph{Absolute continuity})
  $\mu\ll\mu_0$ with density
  $\frac{d\mu}{d\mu_0}(\phi) = Z^{-1}\exp(-V(\phi))$,
  where $V(\phi)=\int_{\mathbb{T}^3}:\!\phi^4\!:(x)\,dx$ is finite $\mu_0$-a.s.

\item[\textbf{(P2)}] (\emph{Positive density})
  The density $\frac{d\mu}{d\mu_0}$ is strictly positive $\mu_0$-a.s.,
  hence $\mu\sim\mu_0$.

\item[\textbf{(P3)}] (\emph{Integrability of shifted interaction})
  For any smooth $\psi:\mathbb{T}^3\to\mathbb{R}$, the functional
  $V(\phi - \psi) = \int_{\mathbb{T}^3}:\!(\phi-\psi)^4\!:(x)\,dx$
  is in $L^p(\mu_0)$ for all $p\ge 1$.
\end{enumerate}

\noindent The main result is:

\begin{theorem}\label{thm:main}
For every smooth function $\psi:\mathbb{T}^3\to\mathbb{R}$ (not necessarily
nonzero), the measures $\mu$ and $T_\psi^*\mu$ are equivalent.
\end{theorem}

%------------------------------------------------------------------
\section{Cameron--Martin Space of the GFF}

The \emph{Cameron--Martin space} of the GFF $\mu_0$ is the Hilbert space
\[
  \mathcal{H} \;=\; H^1(\mathbb{T}^3)
  \;=\; \left\{f\in L^2(\mathbb{T}^3) :
    \|f\|_{H^1}^2 := \langle(-\Delta+m^2)f,f\rangle_{L^2} < \infty\right\}.
\]

\begin{lemma}[Cameron--Martin for GFF]\label{lem:CM}
Let $\psi\in H^1(\mathbb{T}^3)$.  Then $T_\psi^*\mu_0\sim\mu_0$, with
Radon--Nikodym derivative
\[
  \frac{d(T_\psi^*\mu_0)}{d\mu_0}(\phi)
  \;=\; \exp\!\left(\langle(-\Delta+m^2)\psi,\,\phi\rangle
        - \tfrac12\|{\psi}\|_{H^1}^2\right),
\]
where $\langle\cdot,\cdot\rangle$ denotes the $\mu_0$-pairing between
$H^1(\mathbb{T}^3)$ and $\mathcal{D}'(\mathbb{T}^3)$.
\end{lemma}

\begin{proof}
This is the classical Cameron--Martin theorem for Gaussian measures; see e.g.\
Bogachev, \emph{Gaussian Measures}, Chapter~2.  The key point is that
the map $\phi\mapsto\langle(-\Delta+m^2)\psi,\phi\rangle$ is a well-defined
centred Gaussian under $\mu_0$ with variance $\|\psi\|_{H^1}^2$, ensuring
the density is in $L^1(\mu_0)$ and integrates to 1.
\end{proof}

\noindent Since smooth functions are in $H^k(\mathbb{T}^3)$ for all $k$, in
particular $\psi\in H^1(\mathbb{T}^3)$, so Lemma~\ref{lem:CM} applies.

%------------------------------------------------------------------
\section{Wick Powers under Smooth Shifts}

We verify property~(P3) and compute the density of $T_\psi^*\mu$ w.r.t.\ $\mu$.
The key algebraic fact is the following.

\begin{lemma}[Shifted Wick powers]\label{lem:wick}
Let $\psi:\mathbb{T}^3\to\mathbb{R}$ be smooth and let $c(x) = C(x,x)$ be the
(renormalized) variance of the GFF at a point.  Then
\begin{equation}\label{eq:wick-expand}
  :\!(\phi(x)-\psi(x))^4\!:_{\mu_0}
  \;=\;
  :\!\phi(x)^4\!:
  - 4\psi(x)\,:\!\phi(x)^3\!:
  + 6\psi(x)^2\,:\!\phi(x)^2\!:
  - 4\psi(x)^3\,\phi(x)
  + \psi(x)^4,
\end{equation}
where all Wick products $:\!\phi^k\!:$ are taken with respect to $\mu_0$.
\end{lemma}

\begin{proof}
Since $\psi(x)$ is a deterministic smooth function, the Wick ordering
$:\!\cdot\!:_{\mu_0}$ subtracts only the divergent \emph{random} contractions
(i.e.\ contractions involving two $\phi$ factors, not involving $\psi$).
Thus $(\phi(x)-\psi(x))^4$ Wick-ordered w.r.t.\ $\mu_0$ is obtained by
expanding the binomial and Wick-ordering only the $\phi$ factors:
\begin{align*}
  :\!(\phi-\psi)^4\!: \;&=\; :\!\phi^4 - 4\psi\phi^3 + 6\psi^2\phi^2
                         - 4\psi^3\phi + \psi^4\!: \\
  &=\; :\!\phi^4\!: - 4\psi:\!\phi^3\!: + 6\psi^2:\!\phi^2\!:
       - 4\psi^3\,\phi + \psi^4,
\end{align*}
where in the last line we used that Wick ordering is linear, $:\!\psi^k\phi^j\!:
= \psi^k:\!\phi^j\!:$ for deterministic $\psi$, and $:\!\phi\!: = \phi$,
$:\!1\!: = 1$.
\end{proof}

\begin{lemma}\label{lem:integrability}
For every $p\ge1$ and smooth $\psi$, the functional
$\phi\mapsto V(\phi-\psi) = \int_{\mathbb{T}^3}:\!(\phi-\psi)^4\!:(x)\,dx$
belongs to $L^p(\mu_0)$.
\end{lemma}

\begin{proof}
By \eqref{eq:wick-expand},
\[
  V(\phi-\psi) \;=\; V(\phi) - 4\int\psi\,:\!\phi^3\!: + 6\int\psi^2:\!\phi^2\!:
                      - 4\int\psi^3\,\phi + \int\psi^4.
\]
It is a standard result in constructive QFT that $\int_{\mathbb{T}^3}:\!\phi^k\!:(x)\,dx
\in L^p(\mu_0)$ for all $p\ge1$ and $k\le4$ in the $\Phi^4_3$ setting
(see, e.g., Glimm--Jaffe, \emph{Quantum Physics}, or Simon,
\emph{The $P(\phi)_2$ Euclidean Quantum Field Theory}).
Since $\psi\in C^\infty(\mathbb{T}^3)$ is bounded, all correction terms in
\eqref{eq:wick-expand} are bounded multiples of $L^p(\mu_0)$ functions, hence
in $L^p(\mu_0)$.  The term $V(\phi)\in L^p(\mu_0)$ is the $\Phi^4_3$
interaction itself, which is also in $L^p(\mu_0)$.
\end{proof}

%------------------------------------------------------------------
\section{Proof of Equivalence}

\begin{proof}[Proof of Theorem~\ref{thm:main}]
We establish the chain of equivalences
\[
  \mu \;\sim\; \mu_0 \;\sim\; T_\psi^*\mu_0 \;\sim\; T_\psi^*\mu.
\]

\textbf{Step 1: $\mu\sim\mu_0$.}
This is property~(P2): $\mu$ and $\mu_0$ are mutually absolutely continuous
with strictly positive Radon--Nikodym density $\frac{d\mu}{d\mu_0}
= Z^{-1}\exp(-V(\phi)) > 0$ $\mu_0$-a.s.

\textbf{Step 2: $\mu_0\sim T_\psi^*\mu_0$.}
Since $\psi\in C^\infty(\mathbb{T}^3)\subset H^1(\mathbb{T}^3)$, this follows
from the Cameron--Martin Lemma~\ref{lem:CM}.  The Radon--Nikodym derivative
$\frac{d(T_\psi^*\mu_0)}{d\mu_0}$ is strictly positive $\mu_0$-a.s.

\textbf{Step 3: $T_\psi^*\mu_0\sim T_\psi^*\mu$.}
We compute the Radon--Nikodym derivative of $T_\psi^*\mu$ with respect to
$T_\psi^*\mu_0$.  For any measurable set $A\subseteq\mathcal{D}'(\mathbb{T}^3)$:
\[
  (T_\psi^*\mu)(A)
  \;=\; \mu\!\left(T_\psi^{-1}(A)\right)
  \;=\; \mu(A - \psi)
  \;=\; \int_{A-\psi}\frac{d\mu}{d\mu_0}(u)\,d\mu_0(u).
\]
Substituting $\phi = u + \psi$ (i.e., $u = \phi - \psi$):
\[
  (T_\psi^*\mu)(A)
  \;=\; \int_A \frac{d\mu}{d\mu_0}(\phi-\psi)\,d(T_\psi^*\mu_0)(\phi).
\]
Hence
\begin{equation}\label{eq:rn}
  \frac{d(T_\psi^*\mu)}{d(T_\psi^*\mu_0)}(\phi)
  \;=\; \frac{d\mu}{d\mu_0}(\phi-\psi)
  \;=\; Z^{-1}\exp(-V(\phi-\psi)).
\end{equation}
By Lemma~\ref{lem:integrability}, $V(\phi-\psi)\in L^1(T_\psi^*\mu_0)$
(since $T_\psi^*\mu_0\sim\mu_0$ and $V(\phi-\psi)\in L^1(\mu_0)$ by
Lemma~\ref{lem:integrability}), so the density \eqref{eq:rn} is in
$L^1(T_\psi^*\mu_0)$.  Moreover, $Z^{-1}\exp(-V(\phi-\psi)) > 0$
$(T_\psi^*\mu_0)$-a.s., so $T_\psi^*\mu\sim T_\psi^*\mu_0$.

\textbf{Conclusion.}
Combining Steps 1--3:
\[
  T_\psi^*\mu \;\sim\; T_\psi^*\mu_0 \;\sim\; \mu_0 \;\sim\; \mu.
\]
In particular, $\mu$ and $T_\psi^*\mu$ are equivalent.

\medskip
\noindent To make the Radon--Nikodym derivative of $T_\psi^*\mu$ w.r.t.\ $\mu$
explicit, combining Steps 1--3:
\begin{equation}
  \frac{d(T_\psi^*\mu)}{d\mu}(\phi)
  \;=\;
  \frac{Z_\psi}{Z}
  \cdot
  \exp\!\bigl(V(\phi) - V(\phi-\psi)\bigr)
  \cdot
  \exp\!\left(\langle(-\Delta+m^2)\psi,\phi\rangle
    - \tfrac{1}{2}\|\psi\|_{H^1}^2\right),
\end{equation}
where $Z_\psi = \int \exp(-V(\phi-\psi))\,d\mu_0(\phi)$ is the normalising
constant for the shifted measure (finite and positive by
Lemma~\ref{lem:integrability}).  By \eqref{eq:wick-expand},
$V(\phi) - V(\phi-\psi)$ is a polynomial in $\psi$ with coefficients that
are integrable $\Phi^4_3$-Wick monomials, confirming this density is in
$L^1(\mu)$ and positive $\mu$-a.s.
\end{proof}

%------------------------------------------------------------------
\section{Remarks}

\begin{remark}[Mass is not essential]
The proof works for any $m\ge0$ such that the $\Phi^4_3$ measure is
constructed.  The case $m=0$ requires handling the zero mode of the Laplacian
separately, but the essential argument is unchanged.
\end{remark}

\begin{remark}[Non-smooth shifts]
The proof breaks down for $\psi\notin H^1(\mathbb{T}^3)$.  In that case
$T_\psi^*\mu_0\not\sim\mu_0$ by the Cameron--Martin theorem
(translations outside the Cameron--Martin space produce singular measures),
and the conclusion of Theorem~\ref{thm:main} is expected to fail.
\end{remark}

\begin{remark}[Weighted shift]
The same argument applies to the $P(\phi)_2$ measure in 2D, and more
generally to any measure of the form $d\mu\propto\exp(-V(\phi))\,d\mu_0$
where $V$ is a polynomial interaction satisfying integrability in $L^p(\mu_0)$
and $\psi$ lies in the Cameron--Martin space of $\mu_0$.
\end{remark}

%------------------------------------------------------------------
\section*{References}

\begin{itemize}
\item V.~Bogachev, \emph{Gaussian Measures}, AMS, 1998.
\item J.~Glimm and A.~Jaffe, \emph{Quantum Physics: A Functional Integral
  Point of View}, 2nd ed., Springer, 1987.
\item M.~Hairer, ``A theory of regularity structures,''
  \emph{Inventiones Mathematicae}, 198(2):269--504, 2014.
\item B.~Simon, \emph{The $P(\phi)_2$ Euclidean Quantum Field Theory},
  Princeton University Press, 1974.
\end{itemize}

\end{document}
